% II. Проектиране на системата.
\section{Проектиране на системата}
\label{sec:design}

Системата е проектирана около едно ограничение: описанието на услугата трябва да е едно, а
различията между доставчиците изнесени в преброими места. Разпръснат ли се различията из цялото
описание, вече се сравняват две различни системи, а не две среди.

Реализацията се състои от пет изпълними файла върху една обща библиотека:
\code{stratus-worker}, \code{stratus-proxy} (обратен посредник, който разпределя заявките кръгово
и обединява метриките на набора реплики), \code{stratus-loadgen}, \code{stratus-autoscaler} и
\code{stratus-cost}. Работникът не знае при кой доставчик работи: портът, размерът на пула от
нишки и адресът на хранилището идват от променливи на средата. Потокът е следният
\figref{fig:arch}: генераторът изпраща заявки към посредника, посредникът избира реплика по
кръгов ред, всеки работник излага собствени броячи в текстов вид, посредникът ги сумира в един
документ, а контролерът чете документа на равни интервали и решава дали да промени броя реплики.

\begin{figure}[H]
\centering
\begin{tikzpicture}[node distance=0.9cm]
  \node[accent] (lg) {\code{stratus-loadgen}};
  \node[block, right=1.8cm of lg] (px) {\code{stratus-proxy}\\кръгово разпределение\\и обединяване на метрики};
  \node[block, right=1.8cm of px] (w1) {\code{stratus-worker} 1};
  \node[block, below=0.35cm of w1] (w2) {\code{stratus-worker} 2};
  \node[block, below=0.35cm of w2] (wn) {\code{stratus-worker} \textit{N}};
  \node[danger, below=1.4cm of px] (as) {\code{stratus-autoscaler}\\политика за мащабиране};

  \draw[line] (lg) -- node[above,font=\scriptsize] {HTTP} (px);
  \draw[line] (px) -- node[above,font=\scriptsize] {/work} (w1);
  \draw[line] (px) -- (w2);
  \draw[line] (px) -- (wn);
  \draw[line, dashed] (w1.south west) -- node[right,font=\scriptsize,pos=0.35] {/metrics} (px.north east);
  \draw[line] (px.south) -- node[right,font=\scriptsize] {обединени метрики} (as.north);
  \draw[line] (as.east) -- ++(2.6,0) |- node[right,font=\scriptsize,pos=0.25] {брой реплики} (wn.south east);
\end{tikzpicture}
\caption{Структура на системата. Всичко на схемата е общо за двата доставчика}
\label{fig:arch}
\end{figure}

\subsection{Единица работа и метрика за мащабиране}

Сравнението по цена изисква знаменател, който не зависи от средата. За такъв е приета \term{една
единица работа}: едно успешно изпълнение на \code{/work}. Изчислението е функцията за време на
бягство на множеството на Манделброт върху решетка $n \times n$ с таван $m$ итерации над фиксиран
правоъгълник от комплексната равнина, сложност $O(n^2 m)$. Прозорецът е фиксиран нарочно:
преместването му би променило обема работа при същите параметри, без нищо във входа или изхода да
покаже промяната. Стандартната единица е $n = 384$, $m = 500$, при която се изпълняват точно
$18\,690\,064$ итерации, проверени от модулен тест.

Работникът излага брояч \code{stratus\_busy\_seconds\_total} и измерител
\code{stratus\_worker\_threads}, а контролерът изчислява
$\text{натоварване} = \Delta\,\code{busy\_seconds} \,/\, (\Delta t \cdot \sum
\code{worker\_threads})$, тоест както всяка система за наблюдение получава скорост от брояч.
Работникът нарочно не изчислява отношението сам: то би носело неговия прозорец на осредняване, а
не прозореца на контролера. Натоварването на процесора не се използва за сигнал: задачата насища
ядро по построение, така че остава близо до единица и когато наборът се справя, и когато е
задръстен.

\subsection{Политика за мащабиране}

Политиката е следене на цел с две допълнения. Желаният брой реплики е
$r_{\text{жел}} = \lceil r_{\text{тек}} \cdot \text{натоварване} / \text{цел} \rceil$, ограничен
между минимума и максимума. Първото допълнение е \term{зона на нечувствителност}: желание,
различаващо се от текущия брой с по-малко от зададен дял, се приема за шум, иначе контролерът
трепти без край. Второто са \term{несиметрични периоди на изчакване}: увеличаването е позволено
след кратко изчакване, намаляването само след по-дълго, защото закъсняла реакция се плаща от всяка
заявка в опашката, а закъсняло връщане на ресурс се плаща с една реплика.

\subsection{Договор на модулите, идентичност и уговорки}

Двата модула приемат един и същ обект, описващ услугата: име, образ, порт, изчислителен ресурс,
памет, граници на броя реплики, целево натоварване, променливи на средата и етикети. Връщат едни и
същи изходи: входна точка, обектно хранилище, платформа и идентичност. Заради този договор смяната
на доставчика е смяна на стойността на една променлива. Пряко еднакви типове възли
не съществуват, затова съпоставянето е записано преди измерванията: ресурсът се задава в единиците
на AWS Fargate, където $1024$ единици са едно ядро, а модулът за Azure дели на $1024$.

Работникът никога не притежава дълготраен ключ: при AWS това е роля на задачата с права само върху
една кофа, при Azure управлявана идентичност с роля за запис в един профил за съхранение, чиито
ключове за споделен достъп са изключени.

\textbf{Две уговорки.} Първо, натоварването е чисто изчислително: модулите създават хранилището и
подават името и идентичността чрез променливи на средата, но в работника няма клиент за обектно
хранилище, защото обръщението би внесло мрежово чакане в единица работа, чиято цел е да бъде
изцяло изчислителна. Второ, моделът на достъпа е \emph{написан}, но не е \emph{прилаган}: модулите
само следват документираните схеми на ресурсите.
